\newpage
\cleardoublepage
\chapter{Literature Review}\label{chapter:literature_review}

This chapter surveys the literature underlying the work's technical decisions.

\Cref{sec:lit_object_detection} covers object detection, with a focus on the lightweight end of the YOLO family and on the detector-plus-classifier ensemble that is the established alternative for wildlife imagery. \Cref{sec:lit_synthetic_data} covers generative image models, the sim-to-real gap, synthetic-image quality evaluation, and the role of synthetic data in long-tail recognition. \Cref{sec:lit_kd_and_embedded} covers knowledge distillation and the quantization and benchmarking practice that governs whether a trained detector fits an embedded latency budget. \Cref{sec:lit_evaluation_metrics} closes with the evaluation literature this work's own metrics and test-set design draw on.

%%%%%%%%%%%%%%%%%%%%%%%%%%%%%%%%%%%%%%%%%%%%%%%%%
\section{Object Detection for Wildlife Monitoring}\label{sec:lit_object_detection}

Modern object detection descends from two lineages that emerged within a few years of each other. The two-stage, region-proposal lineage first proposes candidate regions and then classifies and refines them, reaching its mature form once the proposal step itself became a learned, end-to-end-trainable network \cite{girshickRichFeatureHierarchies2014a,renFasterRCNNRealTime2017}. A top-down, multi-scale feature hierarchy was later added to it, so that objects of widely differing apparent size are detected by the feature level suited to them \cite{linFeaturePyramidNetworks2017}. That is the normal case in wildlife imagery, where one species occupies a few pixels or most of the frame depending on its distance.

The single-stage lineage discarded the separate proposal step. \textit{YOLO}, for \enquote{you only look once}, reframed detection as a single regression over the whole image \cite{redmonYouOnlyLook2016a}, and multi-scale default boxes narrowed the resulting accuracy gap while retaining one forward pass \cite{liuSSDSingleShot2016a}. Explicitly correcting the foreground-background class imbalance inherent to dense single-stage prediction closed most of what remained \cite{linFocalLossDense2017}.

Two-stage detectors keep a modest accuracy advantage, at the cost of two sequential network evaluations per image. Single-stage detectors trade that accuracy for a fixed, single-pass cost that does not grow with the number of candidate objects in a scene. The target hardware for this work has a fixed, tight per-frame latency budget, so only the single-stage lineage is considered further.

%%%%%%%%%%%%%%%%%%%%%%%%%%%%%%%%%%%%%%%%%%%%%%%%%
\subsection{Challenges in Wildlife Object Detection}\label{sec:lit_wildlife_challenges}

Beyond the general architectural trajectory above, wildlife species detection poses four recurring challenges that motivate specific choices made elsewhere in this work.

First, distinguishing wildlife species is fundamentally a fine-grained visual categorization problem. Unlike COCO-style categories (car, dog, person), many pairs of species differ only in subtle, localized visual cues rather than in gross shape or context \cite{weiFineGrainedImageAnalysis2022}.

This small inter-class variation compounds a second, related problem: several groups of species in this work's target taxonomy are visually near-identical to a general-purpose detector. The plains-zebra/Grévy's-zebra, Eurasian-lynx/bobcat, and small-gazelle groupings, later formalized as confusion clusters in this work's own evaluation methodology (\Cref{sec:evaluation_framework}), are instances of exactly the intra-order visual homogeneity that fine-grained recognition research identifies as its core difficulty \cite{weiFineGrainedImageAnalysis2022}.

Third, ecological image collections are long-tailed by nature: a handful of common, easily-photographed species dominate any naturally-collected corpus, while most species are represented by comparatively few images \cite{cuiClassBalancedLossBased2019,zhangDeepLongTailedLearning2023}. This work's own sourcing effort encountered the same imbalance species by species (\Cref{sec:corpus_composition}).

Fourth, and specific to the wildlife domain, models trained on one visual domain generalize poorly to another. Camera-trap classifiers and detectors that perform well at their training locations degrade sharply at new, unseen camera locations \cite{beeryRecognitionTerraIncognita2018}. Background, lighting, and camera geometry are confounded with location rather than with the animal itself, and this confound is what drives the degradation.

Larger deployed camera-trap classifiers reinforce the same pattern. Cross-project transfer between ecologically similar studies still costs measurable accuracy \cite{williIdentifyingAnimalSpecies2019}, and cross-continental transfer can cost considerably more: one large-scale camera-trap classifier's species-level accuracy fell from $98\%$ in-region to $82\%$ on an out-of-sample Canadian ungulate dataset \cite{tabakMachineLearningClassify2019}. That drop occurred despite the same underlying kind of network having demonstrated near-human accuracy on its own training distribution, achieving $>93.8\%$ species-identification accuracy and rising to $96.6\%$ on the confidently-classified subset \cite{norouzzadehAutomaticallyIdentifyingCounting2018}.

This is precisely the domain gap that motivated rejecting the camera-trap archives of \textit{LILA BC}, a public repository of ecology datasets, as a training-data source for this work (\Cref{sec:sourcing_execution}). Camera traps are static, fixed-position, frequently infrared-triggered sensors, whereas the target deployment is a handheld, daylight, close-to-mid-range binocular observation scenario. The two distributions differ in exactly the location-, lighting-, and viewpoint-confounded ways this line of work documents.

More broadly, the wildlife-monitoring machine-learning community has converged on the view that such domain shift, rather than raw model capacity, is now a primary obstacle to deploying automated recognition at scale \cite{tuiaPerspectivesMachineLearning2022}. Capacity nonetheless sets the outer bound on how many species one deployable detector can separate, which is what the single-stage family below has to supply.

%%%%%%%%%%%%%%%%%%%%%%%%%%%%%%%%%%%%%%%%%%%%%%%%%
\subsection{The YOLO Model Family}\label{sec:lit_yolo_family}

Within the single-stage lineage, the YOLO family evolved along two intertwined axes: from anchor-based to anchor-free box prediction, and from non-maximum-suppression (NMS)-dependent to increasingly NMS-free inference. Anchor boxes were first learned by clustering the training-set box distribution \cite{redmonYOLO9000BetterFaster2017}, then dropped entirely in favor of direct prediction from a decoupled classification and regression head \cite{geYOLOXExceedingYOLO2021}, and the NMS post-processing step was finally removed from inference by consistent dual label assignment during training \cite{wangYOLOv10RealTimeEndtoEnd2024}. Reliably detecting several small objects clustered in one image nonetheless remains unresolved across the generations of the family \cite{nazirYouOnlyLook2023}.

One early design in this lineage is a direct precedent for the taxonomy-aware inference discussed in \Cref{sec:lit_md_speciesnet}: a hierarchical \textit{WordTree} label space that lets a single detector combine detection-only and classification-only training data across a taxonomy \cite{redmonYOLO9000BetterFaster2017}.

Licensing, not architecture, fixes which generation this work can train. \textit{YOLOv5} \cite{ultralyticsComprehensiveGuideUltralytics} is commercially usable only up to commit \texttt{5cdad89}, and later commits require additional licensing that is incompatible with the requirement that the resulting weights be usable in a commercial product. This rules out every newer Ultralytics release independently of any architectural or accuracy comparison, and \textit{YOLOv5s} is accordingly the architecture trained throughout \Cref{sec:augmentation_setups}. None of the family's releases from YOLOv5 onward has a peer-reviewed paper of its own, so citations here are to codebases, release notes, or independent architectural analyses, which is standard practice for this literature.

The nano-scale variants most relevant to embedded deployment sit in the $2.6$--$3.2\,\mathrm{M}$-parameter range (\Cref{sec:target_class_definition}). No paper in the surveyed literature establishes a class-count-versus-accuracy ceiling for backbones of that size. The closest indirect evidence holds class count fixed and varies backbone size instead: a $0.99\,\mathrm{M}$-parameter detector reaches $30.6\%$ average precision on COCO against $25.3\%$ for a $0.91\,\mathrm{M}$-parameter one \cite{yuPPPicoDetBetterRealTime2021,geYOLOXExceedingYOLO2021}. That establishes discriminative capacity as a binding constraint for sub-million-parameter backbones in general, but not how that capacity trades off against a growing number of classes. The $200$--$250$-class ceiling adopted in \Cref{sec:target_class_definition} is therefore this work's own extrapolation from that general capacity-accuracy relationship, not a literature-established figure, and it is treated as such throughout.

A ceiling on what one small detector can resolve makes the established alternative for wildlife imagery worth examining directly.

%%%%%%%%%%%%%%%%%%%%%%%%%%%%%%%%%%%%%%%%%%%%%%%%%
\subsection{The Detector-Plus-Classifier Ensemble and the End-to-End Alternative}\label{sec:lit_md_speciesnet}\label{sec:lit_yolo_vs_ensemble}

A second architecture takes the opposite structural approach: a species-agnostic localizer feeding a separate, much larger species classifier. This work uses it inside its own pipeline as a pseudo-labeling and quality-filtering tool rather than as a deployment target (\Cref{sec:staged_filtering_pipeline,sec:speciesnet_matching}).

\textit{MegaDetector} localizes animals, humans, and vehicles in camera-style imagery without attempting to resolve species identity \cite{beeryEfficientPipelineCamera2019}. Its output crops pass to \textit{SpeciesNet}, a classifier built on an \textit{EfficientNetV2-M} backbone \cite{gadotCropNotCrop2024a} and trained, in its largest reported configuration, on roughly $36\,\mathrm{M}$ camera-trap images spanning more than $2{,}000$ taxonomic labels \cite{gadotCropNotCrop2024}.

Rather than committing to a species-level prediction regardless of confidence, the classification stage applies a confidence-gated taxonomic rollup. When the top species-level prediction falls below a confidence threshold, per-taxon confidence scores are summed upward across the taxonomy until some higher clade (genus, family, and so on) clears it. A low-confidence prediction thereby degrades into a coarser but still informative answer instead of an outright wrong species call \cite{gadotCropNotCrop2024}. A complementary geofencing step rolls up the same way whenever the top prediction falls outside a region-specific allow-list \cite{gadotCropNotCrop2024}.

Two properties make this arrangement the practical accuracy reference for wildlife species recognition from camera-style imagery. Cropping to a detected animal before classifying improves macro-averaged F1 by roughly $25\%$ over whole-image classification on a large, long-tailed camera-trap dataset, and the result reproduces on two further public datasets \cite{gadotCropNotCrop2024}. The rollup then guarantees that a classification too uncertain to name a species still returns a correct genus or family. A single detector trained directly on species-level classes has no such fallback by default. Below its confidence threshold a detection is simply dropped or misclassified, unless an intermediate taxonomic level is deliberately engineered into the label space or the loss, which \textit{WordTree} showed to be possible in principle for a detector \cite{redmonYOLO9000BetterFaster2017}.

The cost is structural, and the two components' separate design does not hide it. A full detector and a large classification backbone both run, and the classifier runs once per detected crop, so total inference time scales with the number of animals in the frame instead of staying constant. \textit{EfficientNetV2-M} on its own is far larger than the nano-scale backbones already established as demanding to fit inside the per-frame latency budget of the \textit{QCS605}'s Hexagon 685 DSP, and that comparison is before counting the detector backbone or the per-crop multiplier.

The choice between the two architectures therefore resolves on deployability rather than on accuracy. The ensemble is more accurate and degrades more gracefully, but fitting it to the target hardware would mean compressing two separately trained networks, each with its own accuracy-versus-size curve. A single end-to-end detector trained on the 225-class taxonomy gives up both the accuracy margin and the taxonomic fallback, but runs at a fixed per-frame cost and leaves one network to compress and export. This work trains the single detector, which moves the accuracy burden onto the training data: a nano-scale detector asked to resolve 225 fine-grained classes must see enough of each one.

%%%%%%%%%%%%%%%%%%%%%%%%%%%%%%%%%%%%%%%%%%%%%%%%%
\section{Synthetic Training Data for Visual Recognition}\label{sec:lit_synthetic_data}

Synthetic imagery enters this work as a remedy for classes whose real-image pool is too small to train on (\Cref{sec:synthetic_data_supplementation}). This section proceeds from the generative models that produce such imagery, to the sim-to-real gap that governs whether it is usable for training at all, to how its quality can be measured, and finally to the specific role it can play for long-tail classes.

%%%%%%%%%%%%%%%%%%%%%%%%%%%%%%%%%%%%%%%%%%%%%%%%%
\subsection{Generative Models for Image Synthesis}\label{sec:lit_generative_models}

Deep generative image modeling was dominated for a decade by adversarial training, in which a generator and a discriminator are optimized against each other in a minimax game \cite{goodfellowGenerativeAdversarialNetworks2014}. The objective is notoriously unstable to optimize, and it admits a degenerate mode in which the generator collapses many inputs onto a single output.

Diffusion models approach synthesis from a different direction. Instead of learning a single feed-forward noise-to-image mapping, they learn to reverse a fixed Markov chain that gradually corrupts an image into noise, trained by predicting the noise added at each step of the forward corruption \cite{hoDenoisingDiffusionProbabilistic2020}. Sample quality surpassed adversarial training without inheriting its instability, at the price of many sequential network evaluations per sample.

Two developments turned this into the text-to-image generation used in this work. \textit{Latent diffusion}, the basis of \textit{Stable Diffusion}, runs the diffusion process in the compressed latent space of a separately pretrained perceptual-compression autoencoder, and introduces cross-attention layers into the denoising network so that arbitrary conditioning modalities, text among them, can be attached \cite{rombachHighResolutionImageSynthesis2022}. This cuts training and inference cost by roughly an order of magnitude against pixel-space diffusion while retaining competitive synthesis quality, and it is what makes locally-hosted image generation practical on a single accelerator at all.

\textit{Classifier-free guidance} governs the trade-off between sample fidelity and sample diversity \cite{hoClassifierFreeDiffusionGuidance2022}. A single network is trained jointly as a conditional and an unconditional diffusion model, by randomly dropping the conditioning signal during training, and at sampling time the two score estimates are linearly extrapolated. Raising the guidance scale pushes samples toward the conditioning text and away from the unconditional prior, sharpening prompt adherence at the expense of variety. Essentially every text-conditioned diffusion model now exposes this as a sampling parameter, and it is one of the generation settings this work has to fix (\Cref{sec:synthetic_generation_pipeline}).

A second axis, orthogonal to the architectural one, is how much natural-language information a model can condition on at all. An open diffusion checkpoint inherits its text-conditioning limit from whichever text encoder it uses. Checkpoints conditioned on \textit{CLIP}, a contrastively trained image and text encoder \cite{radfordLearningTransferableVisual2021}, are capped at $77$ tokens by its transformer text branch. This is a hard architectural ceiling on prompt length, not a soft practical guideline. Later checkpoints pair a CLIP-family encoder with \textit{T5}, an encoder-decoder transformer pretrained at a maximum sequence length of $512$ tokens \cite{raffelExploringLimitsTransfer2023}. That raises usable prompt length by roughly an order of magnitude and measurably improves adherence on complex, highly detailed prompts, while affecting general aesthetic quality only slightly \cite{esserScalingRectifiedFlow2024}.

Large proprietary providers moved in the opposite direction from these openly released, self-hostable checkpoints. Hosted text-to-image and image-editing services, such as \textit{Gemini} image generation and the \textit{GPT-image} family whose lineage traces to \textit{DALL-E} \cite{openaiDALLECreatingImages2021}, publish neither their architectures nor their text-encoder budgets, and in practice accept prompts far longer than any locally-hosted checkpoint can consume.

The gap between a hosted model's effectively unbounded prompt and an open checkpoint's fixed token ceiling is not a peripheral detail for this work. A prompt written for one cannot be reused verbatim for the other without changing what information survives, which is why prompt length is treated as a separately controlled experimental factor in the generator comparison of \Cref{sec:generator_comparison} rather than as a fixed quantity.

%%%%%%%%%%%%%%%%%%%%%%%%%%%%%%%%%%%%%%%%%%%%%%%%%
\subsection{Synthetic Data for Object Detection and the Sim-to-Real Gap}\label{sec:lit_sim_to_real}

Models trained purely on synthetic imagery can transfer to real-world perception tasks, provided the synthetic data is constructed with transfer in mind. This was established in robotics and computer vision long before diffusion-based image generation existed.

\textit{Domain randomization} showed it in the starkest form. An object-localization network trained exclusively on low-fidelity renders, with randomized and deliberately non-photorealistic textures, lighting, camera pose, and distractor objects, reached centimeter-level real-world localization accuracy with no real-image pretraining at all \cite{tobinDomainRandomizationTransferring2017}. The hypothesis behind it is that with sufficient variability in the simulated distribution, the real world appears to the model as one more random variation rather than as a distinct domain, which runs directly counter to the intuition that closing the sim-to-real gap requires photorealism. Extended to full 2D detection, a domain-randomized detector trained on non-photorealistic renders over random real-photo backgrounds matched or beat one trained on a much higher-fidelity synthetic dataset purpose-built to resemble the real \textit{KITTI} test distribution, and fine-tuning it afterward on a modest amount of real data improved results beyond training on real data alone \cite{tremblayTrainingDeepNetworks2018}.

Generative models pose a harder version of the same problem, because the synthetic distribution is now defined implicitly by a learned generator instead of by an explicit, controllable rendering pipeline. Diffusion output layered on top of a strong pretrained feature extractor measurably improves zero-shot classification, but the same synthetic data used to train a classifier from scratch underperforms real data by a wide margin, requiring roughly five times as many synthetic images to match what real images of the same classes achieve \cite{heSyntheticDataGenerative2023}. A strong pretrained representation can partially absorb the gap between generated and real imagery. Raw supervised training cannot.

The gap narrows substantially when the generator itself is fine-tuned on the target domain rather than merely prompted. A diffusion model fine-tuned on the target label space produces images good enough to train a classifier that scores competitively on the real validation set, though still short of real-data training \cite{aziziSyntheticDataDiffusion2023}. Naively prompting an off-the-shelf model with short class-name prompts, by contrast, produces a very poor classifier, and contemporaneous work found that unmodified open-checkpoint output added nothing at all when mixed into real training data \cite{aziziSyntheticDataDiffusion2023}. Raw visual fidelity of the generator is therefore not the operative variable. Alignment between the generator's learned distribution and the target label space and appearance distribution is.

Constraining generation to stay close to a real reference image is a third position on the same axis. Editing real images semantically, rather than generating them from scratch, and using per-class embeddings learned from as few as one labeled example, outperforms both geometric augmentation and unconstrained diffusion augmentation in fine-grained few-shot settings \cite{trabuccoEffectiveDataAugmentation2025}. Keeping the object's own pixels real and synthesizing only the context around it sharpens the caveat. In earlier work by the present author, a detector trained on real car cutouts pasted onto generated backgrounds improved substantially over the unaugmented baseline, but a control filling the same canvas with per-pixel random noise instead of generated context did at least as well at every substitution ratio tested \cite{hedrichGeneratingSyntheticDatasets2026}. The gain followed from the shifted relative-object-size distribution, not from the realism of the surrounding pixels. That study covered one class, one detector, and one evaluation dataset, so it supports only the narrow conclusion that generator realism cannot be assumed to be the operative variable in downstream detector accuracy.

Three failure modes recur across this body of work. The first is a persistent texture and lighting domain shift between generated pixels and real sensor imagery, present even after substantial fine-tuning effort \cite{aziziSyntheticDataDiffusion2023}. The second is a label- and context-alignment gap, in which a generator's default visual associations for a class or scene diverge from the target task's actual distribution unless explicitly corrected \cite{heSyntheticDataGenerative2023,aziziSyntheticDataDiffusion2023}. The third compounds across successive generations rather than within one: \textit{model collapse}, the degeneration of a generative model's learned distribution when it is recursively retrained on data sampled from itself, which first loses the distributional tails and eventually converges on a low-variance, degenerate output \cite{shumailovAIModelsCollapse2024}.

Together these are why a detector in this work is never judged on synthetic images alone (\Cref{sec:evaluation_framework}). A headline metric computed on synthetic imagery alone could look artificially strong precisely because the domain shift and alignment gaps in question are the one thing it cannot probe. Treating synthetic imagery as a supplement whose reliability is actively verified first requires knowing what can be measured about generated imagery at all.

%%%%%%%%%%%%%%%%%%%%%%%%%%%%%%%%%%%%%%%%%%%%%%%%%
\subsection{Evaluating Synthetic Image Quality}\label{sec:lit_synthetic_evaluation}

Generated imagery has historically been assessed numerically through statistics computed in a pretrained image classifier's feature space, rather than by pixel-level comparison against a reference.

The \textit{Inception Score} (IS) is the earliest form of this idea \cite{salimansImprovedTechniquesTraining2016}. Generated images are propagated through an ImageNet-pretrained classifier, and the score rewards a low-entropy per-image label distribution, meaning each image is confidently assigned to one class, while requiring the label distribution across the whole sample to be high-entropy, meaning the sample is diverse. IS never compares a generated sample against any statistic of the real data distribution, and directly optimizing it produces adversarial-example-like artifacts instead of genuine quality gains \cite{salimansImprovedTechniquesTraining2016}.

\textit{Fréchet Inception Distance} (FID) closes exactly that gap \cite{heuselGANsTrainedTwo2018}. Real and generated images pass through the same pretrained activations, each set is modeled as a multivariate Gaussian, and FID is the Fréchet distance between the two Gaussians, $\mathrm{FID} = \|\mu_r - \mu_g\|^2 + \mathrm{Tr}(C_r + C_g - 2(C_r C_g)^{1/2})$. It increases monotonically under independently controlled image disturbances such as noise, blur, implanted occlusion, and cross-dataset contamination, a consistency property IS does not reliably exhibit \cite{heuselGANsTrainedTwo2018}.

\textit{Kernel Inception Distance} (KID) fixes FID's principal statistical weakness \cite{binkowskiDemystifyingMMDGANs2021}. FID's Gaussian-moment estimator is provably biased and needs a large sample before it converges, and that bias is not guaranteed to be consistent across different pairs of distributions, so two FID values computed at the same sample size can rank two generators in the wrong order. KID instead computes an unbiased squared maximum mean discrepancy between the same representations under a polynomial kernel, converging to its asymptotic value even at small sample sizes while keeping FID's sensitivity to the same disturbances. Both are computed for every candidate generator in this work (\Cref{sec:generator_comparison}).

Feature-space metrics do not reliably track human perceptual judgment. Under a structured, blind, multi-rater protocol grounded in psychophysics, human deception rates showed no significant correlation with FID on face generation and only weak, model-dependent correlation on conditional generation \cite{zhouHYPEBenchmarkHuman2019}. The same protocol found that the strongest generator of its generation deceived only marginally more raters than the chance baseline that defines indistinguishability. The claim this work takes from that result is the divergence itself: automatic and human evaluation axes must be triangulated, not substituted for one another.

Realism is also not the only thing that can fail. An image can be photorealistic and still not depict what its prompt asked for, a failure mode that IS, FID, and KID cannot detect at all, since none of them examines the conditioning text. \textit{CLIPScore} measures it directly and without reference captions, by passing image and prompt through CLIP's respective encoders and rescaling their cosine similarity \cite{hesselCLIPScoreReferencefreeEvaluation2021}. Despite requiring no ground truth, it correlates better with human judgments of caption quality than the $n$-gram-overlap metrics it was benchmarked against, which shows that a general-purpose vision-language embedding transfers to quality assessment without task-specific supervision \cite{hesselCLIPScoreReferencefreeEvaluation2021}.

None of these axes measures what ultimately matters when the object is a downstream detector: whether a model trained on the generated imagery performs well on real data. \textit{Classification Accuracy Score} (CAS) operationalizes that train-on-synthetic, test-on-real idea directly \cite{ravuriClassificationAccuracyScore2019}. A classifier is trained exclusively on class-conditional samples from the generative model, then evaluated on the real held-out test set, and its real-test accuracy is the score. Applied to a generator whose IS and FID at the time approached those of real data, CAS exposed a large top-1 accuracy gap against real-data training, and models with markedly worse FID and IS outperformed it. It also isolated individual classes on which the downstream classifier scored zero on the real validation set, a failure a single global FID number hides entirely \cite{ravuriClassificationAccuracyScore2019}.

Three evaluation axes therefore exist and demonstrably disagree: structured blind human rating, automatic distributional and embedding-based proxies, and downstream task utility. This work evaluates its candidate generators along all three and treats none as decisive (\Cref{sec:results_synthetic_data_practicability}), because a generator strong on one need not be strong on another. What that leaves open is which classes synthetic imagery should be spent on in the first place.

%%%%%%%%%%%%%%%%%%%%%%%%%%%%%%%%%%%%%%%%%%%%%%%%%
\subsection{Synthetic Data for Long-Tail and Rare-Class Problems}\label{sec:lit_synthetic_long_tail}

Real-world visual recognition datasets rarely present a uniform class distribution. Remedies divide into class re-balancing, covering re-sampling and cost-sensitive re-weighting, and information augmentation, covering transfer learning and data augmentation \cite{zhangDeepLongTailedLearning2023}. Both families predate the current generation of generative image models, and both were built to make better use of a class's existing samples rather than to manufacture new ones.

Re-sampling attacks the imbalance at the data level. Duplicating minority-class examples verbatim overfits to a narrow decision region, so the foundational method of this family instead synthesizes new minority-class points by interpolating between a sample and its nearest same-class neighbors in feature space \cite{chawlaSMOTESyntheticMinority2002}. Interpolation can only recombine visual information the dataset already holds, the same pose, background, and lighting variation present in the neighbors it interpolates between. It cannot supply a camera angle, habitat, or lighting condition that no real photograph of the class exhibits.

Cost-sensitive re-weighting works on the loss instead of the data. Weighting each class's loss by the inverse of its effective number of samples, a quantity derived from the diminishing marginal information that additional, likely-overlapping samples of a well-represented class contribute, improves long-tailed accuracy across image benchmarks \cite{cuiClassBalancedLossBased2019}. Enforcing a proportionally larger decision-boundary margin for rarer classes, deferred to a later training stage to avoid the optimization difficulties that re-weighting from the start is known to cause, achieves the same from the margin side \cite{caoLearningImbalancedDatasets2019}. Both change how existing images are counted during training, not whether enough of them exist to represent a class's visual diversity.

Re-balancing may not even be the right target. The best feature representations are learned under plain instance-balanced sampling, with no re-balancing during representation learning at all, and long-tailed performance is recovered afterward by re-adjusting the classifier alone \cite{kangDecouplingRepresentationClassifier2020}. That reframes the practical failure of long-tailed recognition as a biased decision boundary rather than a poorly learned representation. It presupposes, however, that a class's representation was learnable from its available images to begin with, and offers no remedy for a class whose real-image pool is too small for any sampling scheme to recover from.

The same limit bounds the augmentation methods surveyed alongside them. Transfer-based and non-transfer variants alike operate in feature space on top of samples the dataset already contains, and none synthesizes new raw images from an external generative prior. The survey is explicit that how best to augment for long-tailed learning remains an open question, and warns separately that class-agnostic augmentation helps head classes as much as tail classes and therefore delivers no relative benefit where it is needed most \cite{zhangDeepLongTailedLearning2023}.

This contrast, rather than any single method, motivates the band structure adopted in this work (\Cref{sec:band_structure}). Re-sampling, cost-sensitive re-weighting, and feature-space augmentation all redistribute information a class's existing image pool already contains, so their ceiling is set by how much visual diversity that pool holds. Image-level generative synthesis is a categorically different intervention, because it manufactures new appearance content, such as camera angle, habitat, lighting, and pose, that a data-poor class's real images cannot supply. Synthetic supplementation is therefore reserved for the classes whose real pool is judged too thin to carry that diversity, and deliberately withheld from better-supplied ones so that they can serve as a real-data control at matched training-set size.

%%%%%%%%%%%%%%%%%%%%%%%%%%%%%%%%%%%%%%%%%%%%%%%%%
\section{Knowledge Distillation and Embedded Inference Feasibility}\label{sec:lit_kd_and_embedded}

Getting a detector trained this way onto the target device raises two further questions: whether a large model's knowledge can be transferred into a small one, and what governs whether the small one meets a latency budget once deployed.

%%%%%%%%%%%%%%%%%%%%%%%%%%%%%%%%%%%%%%%%%%%%%%%%%
\subsection{Knowledge Distillation Fundamentals}\label{sec:lit_knowledge_distillation}

Compressing a model by training a compact one to reproduce a large one's predictions, rather than the original hard labels, predates deep learning's current generation \cite{buciluaModelCompression2006}. Its now-standard form raises the temperature $T$ of the teacher's softmax before the student's training targets are computed \cite{hintonDistillingKnowledgeNeural2015}.

Softening the teacher's output exposes the relative probabilities it assigns to \emph{incorrect} classes, which encode a similarity structure over the label space that a one-hot ground-truth label cannot. An image of a wolf may carry a small probability of being a dog, but that mistake is far more probable than mistaking it for a car. For a 225-class taxonomy of visually confusable mammals, that similarity structure is exactly the information a hard label discards. The student is trained on a weighted combination of two cross-entropy terms: one against the teacher's softened targets, computed at the same high temperature, and one against the ground-truth labels at a temperature of one. The soft-target term's gradient is rescaled by $T^2$, so the two terms remain comparable as $T$ is varied.

Distillation strategies are organized by which part of the teacher supplies the supervision \cite{gouKnowledgeDistillationSurvey2021}. \emph{Response-based} knowledge matches only the teacher's final-layer outputs. \emph{Feature-based} knowledge additionally supervises the student's intermediate representations, in its original form by mapping a chosen intermediate student layer onto a teacher hint layer through a learned regressor, which permits students that are deeper and thinner than the teacher rather than merely smaller \cite{romeroFitNetsHintsThin2015}. \emph{Relation-based} knowledge matches structural relations between samples or feature-map pairs instead of individual predictions.

Transplanting the classification recipe onto a detector is not straightforward. A detector's output combines dense per-location classification and regression predictions under extreme foreground-background imbalance, and distilling teacher and student features uniformly across that imbalance is empirically counterproductive. Distilling foreground and background regions together with equal weighting performs \emph{worse} than distilling either region on its own, because the mismatch between teacher and student attention is concentrated disproportionately in the foreground \cite{yangFocalGlobalKnowledge2022}. Focal and Global Distillation (FGD) instead splits the feature map into foreground and background using the ground-truth boxes, applies teacher-derived spatial and channel attention masks so that the student's limited capacity is spent on the teacher's most salient pixels and channels, and adds a global relational term built from pixel-pair correlations to recover the cross-region context a hard split discards \cite{yangFocalGlobalKnowledge2022}.

The general point matters for interpreting this work's own distillation comparison (\Cref{sec:results_kd_and_embedded}). Plain, spatially unaware logit or feature matching is adequate for image classification, but a detection head needs this kind of foreground/background- and attention-aware treatment before it performs comparably. A naive transplant of the classification recipe should not be expected to match a detection-specific scheme, and a result from one is evidence about that transplant rather than about distillation in general.

A smaller student is only useful if it also runs fast enough on the target device, which is a question about numerical precision, export format, and measurement protocol rather than about training.

%%%%%%%%%%%%%%%%%%%%%%%%%%%%%%%%%%%%%%%%%%%%%%%%%
\subsection{Embedded Inference Feasibility}\label{sec:lit_embedded_inference}

How on-device inference is measured is itself an active concern. An architecture-neutral methodology for comparing inference latency and throughput across heterogeneous hardware and software stacks \cite{reddiMLPerfInferenceBenchmark2020} has been adapted to microcontroller- and DSP-class devices as \textit{MLPerf Tiny}, which reports accuracy, latency, and energy together and standardizes on a common reference inference stack, so that cross-device comparisons are not confounded by inconsistent runtime choices \cite{banburyMLPerfTinyBenchmark2021}. Its measurement protocol deliberately scopes the timed window to model inference alone, explicitly excluding any pre- or post-processing performed around it. This work's own runtime check (\Cref{sec:results_embedded_benchmark}) observes the same methodological control, and additionally excludes one-time effects such as initial model loading and cache and frequency-scaling warm-up.

Export format and runtime are a second variable such a benchmark must hold fixed. The case for a portable, interpreter-based deployment framework over a bespoke per-device compiled pipeline rests on the embedded ecosystem being fragmented across instruction sets, absent operating-system and dynamic-memory support, and vendor-specific kernel optimizations \cite{davidTensorFlowLiteMicro2021}. The same work treats interpreter behavior under multi-threaded and multi-core execution as a first-class design concern rather than an afterthought. Thread and core-affinity choices on a multi-core ARM system-on-chip, such as the \textit{AX Visio}'s, are therefore a further variable to hold fixed across repeated measurements, alongside the export format itself, whether \textit{ONNX}, \textit{TFLite}, \textit{NCNN}, or a vendor-specific format such as Qualcomm's \textit{SNPE} \texttt{.dlc}.

Numerical precision is the cheapest optimization available, and the foundational literature on it targets exactly this work's hardware family. Quantizing a trained network's weights and activations from 32-bit floating point to 8-bit integers after training, known as post-training quantization (PTQ), keeps classification accuracy within $2\%$ of the floating-point baseline across a range of convolutional architectures when weights are quantized per channel and activations per layer \cite{krishnamoorthiQuantizingDeepConvolutional2018}. The resulting speedup is benchmarked at up to $10\times$ over floating-point CPU inference on Qualcomm fixed-point SIMD DSPs of the same class as the \textit{Hexagon 685} \cite{krishnamoorthiQuantizingDeepConvolutional2018}.

Quantization-aware training (QAT) goes further by simulating quantization during training, inserting fake-quantization nodes into the computation graph so that weights and activation ranges adapt to quantization noise before it is applied at inference \cite{jacobQuantizationTrainingNeural2018}. That work frames integer-arithmetic-only hardware such as the Qualcomm Hexagon as its target deployment setting, and reports that QAT narrows the residual accuracy gap to the floating-point baseline further than PTQ alone. The practical distinction is cost against accuracy: PTQ needs only a trained model and a small calibration set, whereas QAT requires a full retraining pass. On the \textit{QCS605} the vendor toolchain implementing both is the Snapdragon Neural Processing Engine (SNPE) together with the AI Model Efficiency Toolkit (AIMET) \cite{QualcommAimet2026}, cited here to identify the toolchain rather than as an independent technical source.

This work surveys quantization for that context but implements neither PTQ nor QAT, nor any other compression technique. The result reported in \Cref{sec:results_embedded_benchmark} is a single runtime measurement of the unoptimized, full-precision model, so the accelerated integer path this literature targets is surveyed here but not exercised there.

Whatever precision the deployed model runs at, its accuracy still has to be reported in a form that a 225-class, look-alike-heavy taxonomy makes legible.

%%%%%%%%%%%%%%%%%%%%%%%%%%%%%%%%%%%%%%%%%%%%%%%%%
\section{Evaluation Metrics for Fine-Grained Wildlife Detection}\label{sec:lit_evaluation_metrics}

Reporting a single mean-average-precision number is sufficient for a generic 80-class benchmark, but is a poor fit for a 225-class, visually confusable, dual-domain wildlife taxonomy. A model that mistakes a plains zebra for a Grevy's zebra has made a fundamentally different error than one that misses the animal entirely. A model evaluated on a synthetic image has also been probed under a different data-generating process than one evaluated on a field photograph.

This section surveys the literature underlying the work's evaluation design. It moves from the general-purpose detection metrics that any object detector is judged against, through the diagnostic instruments developed for understanding \textit{why} a detector fails, to the question of how a test set that mixes real and synthetic imagery should be interpreted at all. The final subsection is the direct literature anchor for the evaluation methodology adopted in \Cref{sec:evaluation_framework}.

\subsection{General Object Detection Metrics}\label{sec:lit_detection_metrics}

The dominant evaluation protocol for object detection was established by the \textit{PASCAL VOC} challenge \cite{everinghamPascalVisualObject2010}. A detection counts as correct if its bounding box overlaps a ground-truth box above a fixed \textit{Intersection-over-Union} (IoU) threshold, the ratio of the intersection of the two boxes to their union. That threshold is conventionally set to $0.5$, a deliberately low value chosen to absorb the inherent subjectivity of drawing a tight box around a non-convex object such as an animal with limbs extended \cite{everinghamPascalVisualObject2010}.

Every prediction above the confidence threshold is then a true or a false positive, from which precision, the fraction of predictions that are correct, and recall, the fraction of ground-truth objects that are detected, follow as functions of that threshold. \textit{Average Precision} (AP) summarizes the resulting precision-recall curve as the mean interpolated precision at a fixed set of recall levels, and mean Average Precision (mAP) averages AP across classes \cite{everinghamPascalVisualObject2010}. This formulation was adopted because a simple accuracy count is misleading under the heavy imbalance between the small number of true objects and the very large number of candidate regions a detector must reject, and it has remained the basis of essentially every detection benchmark since \cite{everinghamPascalVisualObject2010,zouObjectDetection202023}.

Two refinements produced the metric vector used throughout modern detection research. A fixed IoU of $0.5$ rewards coarse localization, because it cannot distinguish a tightly fitted box from a loose one as long as both clear the bar. Averaging AP over ten IoU thresholds spanning $0.5$ to $0.95$ in steps of $0.05$, denoted AP@[.5:.95], as the \textit{COCO} benchmark does, rewards progressively tighter localization instead \cite{linMicrosoftCOCOCommon2015,zouObjectDetection202023}. The single-threshold AP50 and the stricter AP75 are retained beside it as complementary operating points, the first closer to the original PASCAL protocol and the second probing high-precision localization specifically \cite{zouObjectDetection202023}. Separately, AP credits a detector only for the boxes it does emit, so it says nothing about ground-truth objects that are never proposed at all. \textit{Average Recall} (AR) measures that upper bound on achievable recall independently of the confidence ranking, reported at a fixed budget of the top $1$, $10$, or $100$ detections per image \cite{zouObjectDetection202023}.

The resulting twelve-statistic vector covers AP@[.5:.95], AP50, AP75, AP split by object size, and AR at the three detection budgets split the same way. This work retains it unchanged as its own baseline reporting format (\Cref{sec:evaluation_framework}), because it is the bridge that keeps the work's numbers legible against published results. It is nonetheless too coarse on its own to answer the questions a 225-class, look-alike-heavy wildlife taxonomy raises.

\subsection{Hierarchical Evaluation and Error-Type Decomposition}\label{sec:lit_hierarchical_evaluation}\label{sec:lit_error_decomposition}

A single flat AP computed over all 225 classes cannot distinguish two very different errors: finding an animal and naming the wrong species inside a visually confusable group, such as a plains zebra scored as a Grevy's zebra, and naming a completely unrelated class, such as a wolf scored as a squirrel. Nor does AP say \textit{why} a detector underperforms, whether it fails to localize objects at all, localizes them but assigns the wrong label, or duplicates detections on the same object. From a practical-deployment standpoint these are not equivalent failures, and they call for different fixes.

\textcite{hoiemDiagnosingErrorObject2012} supply the precedent for both distinctions. Every false positive is partitioned into one of four mutually exclusive categories: localization error, meaning a detection that overlaps a true object of the correct class with insufficient IoU, which they treat as including duplicate detections on an already-detected object, confusion with a semantically similar object, confusion with a dissimilar labeled object, and confusion with unlabeled background. Similarity is defined by small sets of semantically related categories, such as all vehicles or all animals, fixed \textit{a priori} from domain knowledge about which categories are related, before any detector is run, and never read off a trained model's confusion matrix afterward.

Two findings from that decomposition carry over to this work. Within-group confusion and pure localization error together account for the large majority of high-confidence false positives, markedly more than confusion with dissimilar or unlabeled background clutter. And reporting the AP improvement that would follow from eliminating each category in isolation identifies which failure mode is actually worth addressing, a conclusion invisible in a single aggregate number. Within-group confusion is precisely the failure mode this work expects to dominate among the zebra, elephant, lynx-caracal, and gazelle-antelope look-alike clusters of the target taxonomy.

The a-priori requirement is the load-bearing part for this work's own grouping strategy (\Cref{sec:evaluation_framework}). A taxonomic or curated visual-similarity rollup, decided independently of any candidate model's behavior, is admissible. A grouping mined from a model's own confusion matrix is circular, because it would define the very categories that model is then evaluated against.

The literature surveyed for this work contains no paper performing such a rollup at species, genus, and family resolution specifically, and none that formally compares a taxonomy-driven grouping against an empirical-confusion-driven one as competing evaluation designs. The closest available anchor is the general a-priori-grouping principle above, which this work extends from a handful of hand-picked semantic sets to a full taxonomic hierarchy. The four-way decomposition is adopted alongside it as a corroborating instrument, used to check whether an accuracy loss between the coarse and fine granularity levels is attributable to classification confusion among look-alike species, as expected, or to a hidden localization or duplicate-detection problem. Modern off-the-shelf toolboxes extend the same partition with a missed-ground-truth category, but the survey conducted for this work could not establish a verifiable citation for the specific implementation used in \Cref{chapter:results}, so the method is anchored here in its original source rather than in a named downstream implementation.

Both instruments still presuppose a single, homogeneous test set.

\subsection{Evaluating Across a Real/Synthetic Test-Domain Split}\label{sec:lit_real_synthetic_evaluation}

Every evaluation protocol discussed so far is computed over one homogeneous test set of real photographs. This work's dataset instead ships two test domains: an uneven, per-class real test set, and a class-balanced $225\times50$ synthetic test set (\Cref{sec:synthetic_data_supplementation}).

Standard practice when synthetic imagery is involved in training is to keep evaluation strictly on real, established test sets, rather than blending synthetic images into the test condition itself. Domain-randomization work tackling the same sim-to-real problem this work faces evaluates its detector exclusively on the real-world \textit{KITTI} test set, and never mixes synthetic renders into the test condition, even though a synthetic dataset rendered to match that test domain was available as a comparison point \cite{tremblayTrainingDeepNetworks2018}. Work evaluating whether modern generative-model output is usable as training data likewise benchmarks its synthetically trained models against standard, unmodified real test sets rather than any blended variant, and attributes the resulting gaps to a domain gap between the generated and real image distributions \cite{heSyntheticDataGenerative2023}.

The general problem of a mismatch between the distribution a model is evaluated on and the one it will actually encounter is the subject of a substantial domain-generalization literature \cite{zhouDomainGeneralizationSurvey2023}, tracing back to early work establishing that benchmark datasets themselves carry systematic biases that limit how their measured performance transfers \cite{UnbiasedLookDataset}. A synthetic-to-real train/test shift is one instance of it.

None of this literature, however, treats a synthetic partition as a component to be folded into a single headline test metric alongside real data. Where synthetic data appears in a test protocol at all, it is either the sole test domain, isolating a specific capability, or is kept rigidly separate from the real test condition, precisely so that the domain-shift gap stays measurable.

The surveyed literature therefore offers \textbf{no established precedent} for two specific choices this work makes: reporting a headline detection metric computed over the union of a real and a synthetic test set, and using a class-balanced synthetic test set as a diagnostic instrument for classes where the real test data is thin or noisy. Both are argued for on their own terms in \Cref{sec:evaluation_framework}, which also defines the real-only breakout reported alongside the mixed figure and the real-versus-synthetic watchdog that would force a revision of the evaluation axes. The literature above is context for that argument, showing what the established alternative looks like and why departing from it needs its own justification. It is not literature that itself endorses a mixed real-and-synthetic headline metric.
